\SourceLeaf{210}
äussersten Werthe von \(u\) gleich gross wird, wodurch es gegen die Mitte zu seinen kleinsten Werth erhält. Sind diese äussersten Werthe von \(u\) diese \(u^0\) und \(u'\), so wird man demnach setzen müssen:
\[
\lambda=\frac{\log\sin u'-\log\sin u^0}{\log\mathrm{tang}\frac{1}{2}u'-\log\mathrm{tang}\frac{1}{2}u^0}
\]
\noindent Die Blätter von Herrn Professor \textsc{Harding}'s Sternkarten Nr. 19---26 sind nach dieser Projection gezeichnet.

\Article{11}
Man kann die allgemeine Auflösung für das im vorhergehenden Artikel behandelte Beispiel noch in einer andern Form aufstellen, die wir ihrer Eleganz wegen hier noch beifügen zu müssen glauben.

In Folge des im 6. Art. Vorgetragenen wird, da
\[
\mathrm{tang}\tfrac{1}{2}u(\cos t+i\sin t)
\]
\noindent eine Function von
\[
t+i\log\mathrm{cotang}\tfrac{1}{2}u
\]
\noindent ist, und
\[
\mathrm{tang}\tfrac{1}{2}u(\cos t+i\sin t)=\frac{\sin u\cos t+i\sin u\sin t}{1+\cos u}=\frac{x+iy}{a+z}
\]
\noindent die allgemeine Auflösung auch durch
\[
X+iY=f\frac{x+iy}{a+z},\quad X-iY=f'\frac{x-iy}{a+z}
\]
\noindent dargestellt werden können, d. i. \(X\) muss dem reellen und \(iY\) dem imaginären Theil von \(f\frac{x+iy}{a+z}\) gleich gesetzt werden, indem \(f\) eine willkürliche Function bezeichnet. Anstatt \(f\frac{x+iy}{a+z}\) kann man, wie man leicht sieht, auch eine willkürliche Function von \(\frac{y+iz}{a+x}\), oder von \(\frac{z+ix}{a+y}\) nehmen.

\Article{12}
Wir wollen viertens die Darstellung der Oberfläche des Revolutions-Ellipsoids in der Ebne betrachten. Es seien \(a\) und \(b\) die beiden halben Hauptaxen des Ellipsoids, so dass
\[
\begin{aligned}
x&=a\cos t\sin u\\
y&=a\sin t\sin u\\
z&=b\cos u
\end{aligned}
\]
